\section{Experiments}\label{section::Experiments}


% -------------------------以下是中文版本---------------------------------------%
本章围绕“检测—跟踪—计数”的整体流程，系统评估了所提出的人群头部检测与站台场景多目标跟踪计数方法。我们以YOLO11m作为检测器，结合基于MobileNet-B0 的 ReID 外观描述模型，并在 DeepSORT框架内分别接入三类运动模型以开展实验，包括：8维度匀速基准模型（标准DeepSORT卡尔曼滤波器）；12维度加速度模型以及结合针孔相机集合约束的5维度物理模型。在Data章节所提供的数据上进行评估，其中检测器训练与微调所用的静态图像数据集与多段视频序列互相补充，既涵盖了通用的拥挤场景，也覆盖了列车进站期间不同密度的平台场景。

实验设置

1.检测器YOLO1m训练与微调：第一阶段使用CrowdHuman数据集训练YOLO11m，第二阶段使用我们整合并标注的 Open Sensor Data for Rail 2023、RailEye3D 与自建 TrainPlatformCrowd 三个来源的共2,000 张头部标注图像进行微调，得到 YOLO11freeze10.pt。
如图yolopre2_Results Summary，展示了第一阶段训练的结果，如表yolopre2_results所示，我们训练了100epochs，在CrowdHuman数据集上达到了79.4\%的mAP(50),以及54.7%的mAP(50-95)
\begin{table}[h!]
    \centering
    \setlength{\tabcolsep}{12pt} % 调整列间距
    \renewcommand{\arraystretch}{1.3} % 调整行高
    \begin{tabular}{|c|c|c|c|c|c|c|}
    \hline
    \textbf{Class} & \textbf{Images} & \textbf{Instances} & \textbf{Box(P)} & \textbf{R} & \textbf{mAP50} & \textbf{mAP50--95} \\
    \hline
    \rowcolor{blue!20}
    all & 4370 & 98568 & 0.862 & 0.715 & 0.794 & 0.547 \\
    \hline
    \end{tabular}
    \end{table}
    \caption{YOLO11m在CrowdHuman数据集上的训练结果}
    \label{tab:yolopre2_results}
如图，yolofreeze10_Results Summary，展示了第二阶段训练的结果，如表yolofreeze10_results所示，我们训练了100epochs，在Open Sensor Data for Rail 2023、RailEye3D 与自建 TrainPlatformCrowd 三个来源的共2,000 张头部标注图像上达到了98.0\%的mAP(50),以及81.6%的mAP(50-95)

\begin{table}[h!]
    \centering
    \setlength{\tabcolsep}{12pt} % 调整列间距
    \renewcommand{\arraystretch}{1.3} % 调整行高
    \begin{tabular}{|c|c|c|c|c|c|c|}
    \hline
    \textbf{Class} & \textbf{Images} & \textbf{Instances} & \textbf{Box(P)} & \textbf{R} & \textbf{mAP50} & \textbf{mAP50--95} \\
    \hline
    \rowcolor{blue!20}
    all & 1676 & 9352 & 0.965 & 0.946 & 0.98 & 0.816 \\
    \hline
    \end{tabular}
    \end{table}
    \caption{YOLO11m在CrowdHuman数据集上的训练结果}
    \label{tab:yolofreeze10_results}

按 \S\ref{section::Methods} 所述，我们选用YOLO11m作为精度—速度折中版本。第一阶段在CrowdHuman数据集中训练获得通用头部检测能力，第二阶段以 \texttt{backbone=10} 冻结微调策略，得到YOLO11freeze10.pt作为最终检测器权重。跟踪阶段默认使用高置信度阈值 \(\geq 0.7\) 与 IOU阈值为0.5，旨在减少非人头目标进入候选和匹配的，提高跟踪速度和准确性，获得更好的评估结果。

检测器消融实验
为了探究微调后的检测器在不同设置下的性能，我们进行了消融对比实验。首先，使用相同的设备，相同的环境，在400张验证集图像上测试不同超参数组合这样选择数据集的好处是节省了资源，无需额外的数据收集和标注，且使用同一组数据集保证了实验的一致性，最重要的是直接在目标场景上进行了验证，也符合一般的做法。
我们先同一尺寸，使用常用的图像输入尺寸imagesize=640，探究统一尺寸下的不同置信度和不同IoU阈值的性能差异，然后再根据实验结果，设置不同的输入尺寸，探究不同的输入尺寸对性能的影响。
如表所示，IOU和置信度的实验超参数设置如表所示：
表
        - 置信度阈值：0.3, 0.5, 0.7, 0.9
        - IoU阈值：0.3, 0.5, 0.7, 0.9  

经过试验后，结果如下表所示P25_YOLO11m_Ablation，可以看到，当置信度阈值为0.7，IoU阈值为0.5，检测器的性能最好。
然后，我们根据上面的实验结果，设置Iou=0.5，置信度阈值为0.7，探究不同的输入尺寸对性能的影响。结果如下表所示，
- 输入尺寸：640, 736, 832, 960



2.ReID（MobileNet）实验设置与结果
根据第4章的方法，来自Zhu等人的研究\cite{218}，考虑效率和准确度，如图P26_EfficientNet-B0_Results.png所示，我们采用EfficientNet-B0作为骨干网络，重新训练ReID模型来适合我们系统。
训练ReID的自建数据集：使用10个视频序列构建训练数据集，包含238个有效跟踪目标，如表ReID_Data_Distribution所示用来训练ReID模型的数据集。根据表格的数据，我们的目标接近正方形，如图P24_aspect_ratio_distribution.png和P24_1_width_vs_height.png显示，因此我们采取正方形尺寸，然后经过智能缩放与填充预处理。
分析检测框的尺寸的分析图，
四种不同输入尺寸的消融实验 64,96,128,256
考虑到不同输入尺寸训练出来的模型权重不同，模型学习到的特征不同，对细节信息的捕获能力也不同。我们设计了对比实验，对比不同输入尺寸对模型速度和精度的影响。结果如表所示。


基于对84,354个检测框的统计分析，128×128被确立为在计算效率与特征丰富度之间达到最优平衡的输入尺寸。该尺寸能够覆盖数据集中99.5%以上的检测框，即使最大目标（259.9×289.0）在保持长宽比的智能缩放策略下仍可有效表示，同时为面积小于500px²的小目标（占比13.3%）提供约100倍的像素增益，显著增强其特征表达能力。与64×64相比，128×128显著减轻了中大型目标（如面积超过2000px²，占比34.8%）的细节损失；与256×256相比，在仅0.5%的目标面积超过20,000px²的情况下，其计算与内存开销更优——内存占用降低75%，更适合实时推理与资源受限场景。此外，数据集检测框的平均长宽比为0.920（标准差0.166），接近正方形，使128×128的正方形输入可最大限度减少填充区域。综上，128×128在适配绝大多数目标尺寸分布、保持较高像素密度、提升训练并行性与推理速度方面表现出综合优势，是ReID模型推荐的输入尺寸。


3.三类运动模型

为验证运动建模对平台场景多目标跟踪稳定性与ID一致性的影响，我们在相同检测与 ReID 设置下，对比以下三种卡尔曼滤波方案（接口均与 DeepSORT 兼容），我们使用MOT-TPCH数据集进行评估，该数据集包含20个视频序列，如图P20_MOT-TPCH_Evaluation20所示，该20个视频用来评估3个跟踪算法的MOT评估指标及其计数指标。

在分析之前，我们尝试对所跟踪的目标进行自然语言建模，如图，以为例的真实标签的轨迹图，我们知道，在2D层面上，这是一个以匀加速的一个线性运动过程。

以单独一个目标为例，该目标在远处以较小的尺寸被识别，然后随着火车的前进，目标逐渐变大，直到完全进入视野，然后再逐渐退出画面。对于火车刚进入站台时，目标相对火车的速度让我们发现几乎所有的目标可视静止不动。
% ---------------------------------------% 8维匀速基准模型（CV-8D）------------------------------------------------------------
\paragraph{8维匀速基准模型（CV-8D）。} 该模型假设常速度运动，使用标准线性观测。其实现对应仓库中的 \texttt{deep\_sort/kalman\_filter.py} 的基类实现，作为广泛采用的基准。其中，

\textbf{状态向量为：}
\begin{equation}
\mathbf{x}_k = [x, y, a, h, \dot{x}, \dot{y}, \dot{a}, \dot{h}]^T \in \mathbb{R}^{8}
\end{equation}

\textbf{状态转移矩阵为：}
\begin{equation}
\mathbf{F} = \begin{bmatrix}
1 & 0 & 0 & 0 & \Delta t & 0 & 0 & 0 \\
0 & 1 & 0 & 0 & 0 & \Delta t & 0 & 0 \\
0 & 0 & 1 & 0 & 0 & 0 & \Delta t & 0 \\
0 & 0 & 0 & 1 & 0 & 0 & 0 & \Delta t \\
0 & 0 & 0 & 0 & 1 & 0 & 0 & 0 \\
0 & 0 & 0 & 0 & 0 & 1 & 0 & 0 \\
0 & 0 & 0 & 0 & 0 & 0 & 1 & 0 \\
0 & 0 & 0 & 0 & 0 & 0 & 0 & 1
\end{bmatrix}_{8 \times 8}
\end{equation}

\textbf{过程噪声协方差矩阵为：}
\begin{equation}
\mathbf{Q}_k = \text{diag}(\sigma_{pos}^2 h, \sigma_{pos}^2 h, 10^{-4}, \sigma_{pos}^2 h, \sigma_{vel}^2 h, \sigma_{vel}^2 h, 10^{-10}, \sigma_{vel}^2 h)
\end{equation}
其中 $\sigma_{pos} = 1/20$，$\sigma_{vel} = 1/160$。

\textbf{观测向量为：}
\begin{equation}
\mathbf{z}_k = [x, y, a, h]^T \in \mathbb{R}^{4}
\end{equation}

\textbf{观测矩阵为：}
\begin{equation}
\mathbf{H} = \begin{bmatrix}
1 & 0 & 0 & 0 & 0 & 0 & 0 & 0 \\
0 & 1 & 0 & 0 & 0 & 0 & 0 & 0 \\
0 & 0 & 1 & 0 & 0 & 0 & 0 & 0 \\
0 & 0 & 0 & 1 & 0 & 0 & 0 & 0
\end{bmatrix}_{4 \times 8}
\end{equation}
我们使用CV-8D作为基准模型，在MOT-TPCH数据集上针对20个视频序列进行评估，如表P21_CV-8D_Results所示，我们使用CV-8D模型，在MOT-TPCH数据集上达到了23\%的MOTA，900\%的IDF1，90.0\%的IDP，90.0\%的IDR。
如图所示P21_CV-8D_Results_videoFrame，是CV-8D模型的评估视频过程中的某一帧，其中展示了当前统计的左右各计数带的人数，以及当前的计数结果，以及目标被跟踪的情况。


% ---------------------------------------% 12维加速度模型（CA-12D）------------------------------------------------------------
\paragraph{12维加速度模型（CA-12D）。} 在基准模型上扩展加速度分量，状态为 \([x, y, a, h, v_x, v_y, v_a, v_h, a_x, a_y, a_a, a_h]^\top\)，测量仍为 \([x, y, a, h]^\top\)。相较 CV-8D，能更好刻画进站减速、短时加减速与遮挡恢复阶段的状态变化。实现参考 \texttt{deep\_sort/kalman\_filter.py} 中的 \texttt{AccelerationKalmanFilter}（提供速度/加速度平滑与缺帧保守预测）。


\textbf{状态向量为：}
\begin{equation}
\mathbf{x}_k = [x, y, a, h, \dot{x}, \dot{y}, \dot{a}, \dot{h}, \ddot{x}, \ddot{y}, \ddot{a}, \ddot{h}]^T \in \mathbb{R}^{12}
\end{equation}

\textbf{状态转移矩阵为：}
\begin{equation}
\mathbf{F} = \begin{bmatrix}
\mathbf{I}_4 & \Delta t \cdot \mathbf{I}_4 & \frac{1}{2}\Delta t^2 \cdot \mathbf{I}_4 \\
\mathbf{0}_{4 \times 4} & \mathbf{I}_4 & \Delta t \cdot \mathbf{I}_4 \\
\mathbf{0}_{4 \times 4} & \mathbf{0}_{4 \times 4} & \mathbf{I}_4
\end{bmatrix}_{12 \times 12}
\end{equation}

\textbf{运动方程为：}
\begin{align}
\mathbf{p}_{k+1} &= \mathbf{p}_k + \mathbf{v}_k \Delta t + \frac{1}{2}\mathbf{a}_k \Delta t^2 \\
\mathbf{v}_{k+1} &= \alpha_v \mathbf{v}_{k+1}^{pred} + (1-\alpha_v) \mathbf{v}_k, \quad \alpha_v = 0.85 \\
\mathbf{a}_{k+1} &= \alpha_a \mathbf{a}_{k+1}^{pred} + (1-\alpha_a) \mathbf{a}_k, \quad \alpha_a = 0.7
\end{align}

\textbf{过程噪声协方差矩阵为：}
\begin{equation}
\mathbf{Q}_k = \text{diag}(\sigma_{pos}^2 h, \sigma_{vel}^2 h, \sigma_{acc}^2 h) \otimes \mathbf{I}_4
\end{equation}
其中 $\sigma_{pos} = 1/20$，$\sigma_{vel} = 1/160$，$\sigma_{acc} = 1/10$。

\textbf{观测向量为：}
\begin{equation}
\mathbf{z}_k = [x, y, a, h]^T \in \mathbb{R}^{4}
\end{equation}

\textbf{观测矩阵为：}
\begin{equation}
\mathbf{H} = [\mathbf{I}_4, \mathbf{0}_{4 \times 4}, \mathbf{0}_{4 \times 4}]_{4 \times 12}
\end{equation}


我们使用CA-12D作为基准模型，在MOT-TPCH数据集上针对20个视频序列进行评估，如表P22_CA-12D_Results所示，我们使用CA-12D模型，在MOT-TPCH数据集上达到了23\%的MOTA，900\%的IDF1，90.0\%的IDP，90.0\%的IDR。
如图所示P22_CA-12D_Results_videoFrame，是CA-12D模型的评估视频过程中的某一帧，其中展示了当前统计的左右各计数带的人数，以及当前的计数结果，以及目标被跟踪的情况。


% ---------------------------------------% 5维物理模型（Phys-5D）------------------------------------------------------------
\paragraph{5维物理模型（Phys-5D）。} 融合针孔相机几何约束，将像素高度 \(h\) 与深度 \(Z\) 关联：\(Z = fH/h\)。状态取 \([u, v, h, Z, \dot Z]^\top\)，测量为 \([u, v, h]^\top\)。通过 \((X,Y,Z) \leftrightarrow (u,v,h)\) 的投影/反投影在预测时显式更新 \(Z\) 与 \(h\)，并对 \(\dot Z\) 引入减速度约束以贴合列车相对运动的物理先验。实现参考 \texttt{physical3dkf\_5d.py} 与 \texttt{deep\_sort/custom\_kf5d.py}。

\textbf{状态向量为：}
\begin{equation}
\mathbf{x}_k = [u, v, h_{px}, Z, \dot{Z}]^T \in \mathbb{R}^{5}
\end{equation}

\textbf{几何约束关系：}
\begin{equation}
h_{px} = \frac{f \cdot H}{Z}, \quad Z = \frac{f \cdot H}{h_{px}}
\end{equation}
其中 $f$ 为相机焦距（像素），$H = 0.3m$ 为行人头高。

\textbf{深度动力学模型：}
\begin{align}
Z_{k+1} &= Z_k + \dot{Z}_k \Delta t + \frac{1}{2}a_Z \Delta t^2 \\
\dot{Z}_{k+1} &= \dot{Z}_k + a_Z \Delta t
\end{align}

\textbf{自适应减速度约束：}
\begin{equation}
a_Z = -\text{sign}(\dot{Z}) \cdot \min\left(a_{max}, \frac{\dot{Z}^2}{2Z}\right), \quad a_{max} = 1.0 \text{m/s}^2
\end{equation}

\textbf{相机内参配置：}
根据视频分辨率和裁剪区域：
\begin{itemize}
    \item 2K右侧：$f = 795.0$, $(c_x, c_y) = (268.0, 137.0)$
    \item 2K左侧：$f = 795.0$, $(c_x, c_y) = (1268.0, 137.0)$
    \item 4K右侧：$f = 1600.3$, $(c_x, c_y) = (515.0, 153.0)$
    \item 4K左侧：$f = 1600.3$, $(c_x, c_y) = (1789.0, 153.0)$
\end{itemize}

\textbf{过程噪声协方差矩阵为：}
\begin{equation}
\mathbf{Q}_k = \text{diag}(\sigma_{pos}^2 h_{px}, \sigma_{pos}^2 h_{px}, \sigma_{pos}^2 h_{px}, \sigma_{vel}^2 h_{px}, \sigma_{vel}^2 h_{px})
\end{equation}
其中 $\sigma_{pos} = 1/2$，$\sigma_{vel} = 1/15$。

\textbf{观测向量为：}
\begin{equation}
\mathbf{z}_k = [u, v, h_{px}]^T \in \mathbb{R}^{3}
\end{equation}

\textbf{观测矩阵为：}
\begin{equation}
\mathbf{H} = \begin{bmatrix}
1 & 0 & 0 & 0 & 0 \\
0 & 1 & 0 & 0 & 0 \\
0 & 0 & 1 & 0 & 0
\end{bmatrix}_{3 \times 5}
\end{equation}

我们使用Phys-5D作为基准模型，在MOT-TPCH数据集上针对20个视频序列进行评估，如表P23_Phys-5D_Results所示，我们使用Phys-5D模型，在MOT-TPCH数据集上达到了67.21\%的MOTA，76.29\%的IDF1，82.23\%的IDP，71.51.0\%的IDR。
如图所示P23_Phys-5D_Results_videoFrame，是Phys-5D模型的评估视频过程中的某一帧，其中展示了当前统计的左右各计数带的人数，以及当前的计数结果，以及目标被跟踪的情况。


4.跟踪与计数评估：在自建的数据集MOT-TPCH上进行评估。该数据集按\ref{section::data}所述经由时间/空间裁剪与去音频处理，分辨率介于 1920\(\times\)1080 与 2304\(\times\)1296/1536\(\times\)864 之间，并标注连续头部轨迹（MOT 格式）。
其中20个视频序列用于评估，如表P20_MOT-TPCH_Evaluation20所示，该20个视频用来评估3个跟踪算法的MOT评估指标及其计数指标。

4.1 计数方法的消融实验
根据上面的结果，我们采取虚拟计数带的方法进行计数，并且设置计数带在距离图像边界0.2个宽度的位置，该位置的人员跟踪效果最稳定，效果最好。

传统的计数线的方法：单线计数方法存在一定的局限性，比如当目标被遮挡，计数线会遗落暂时被遮挡的目标，比如在计数线上有3个人，但是前面两个人挡住了第三个人，计数线只计数了两个值；当短暂的位移后，如果是计数带，则会准确地计数到3个目标。我们设置计数线在距离图像边界0.2个宽度的位置，该位置的人员跟踪效果最稳定，效果最好。

虚拟计数带的方法：减少因短暂遮挡导致的漏计数，提高计数稳定性，并且最后帧补偿机制提升完整性，解决了最后画面目标未通过计数线或者计数带的漏记问题。如图，我们设置计数带的位置，并进行计数，可以看到计数带的方法能够准确地计数到3个目标。

两种计数方法的结果比较。

在确定使用计数带进行计数后，我们进行了如下消融实验，对比了不同的计数带宽度的计数性能，结果如下表所示P24_AreaCounting_Ablation，对比了计数带和计数线方法在：


